\chapter{Hyperbolic Dehn surgery}

\section{Ideal tetrahedra and gluing}

\begin{definition}[Ideal tetrahedron shape]
\label{thurston-ch04-ideal-tetrahedron-shape}
\group{hyperbolic geometry}
\source{thurston:page-0056,thurston:page-0058,thurston:page-0059}
\notready
For an oriented ideal tetrahedron, the link of each ideal vertex is a Euclidean triangle up to oriented similarity. If the vertices of such a triangle are labelled cyclically and \(v\) is placed at \(0\), \(u\) at \(1\), its shape parameter is \(z(v)=t\in\C\) with \(\operatorname{Im}t>0\). The other vertex parameters are \(z(t)=(t-1)/t\) and \(z(u)=1/(1-t)\). The three edge parameters are constant on opposite edges.
\end{definition}

\begin{proposition}[Ideal tetrahedron shape identities]
\label{thurston-ch04-shape-identities}
\group{hyperbolic geometry}
\source{thurston:page-0057,thurston:page-0059}
If \(z_1,z_2,z_3\) are the three distinct shape parameters in cyclic order, then
\[
 z_1z_2z_3=-1,\qquad 1-z_1+z_1z_2=0.
\]
Opposite edges have equal dihedral angles and equal shape parameters, and the three angles at every ideal vertex sum to \(\pi\).
\notready
\end{proposition}

\begin{proof}
\sketch The two alternate vertex normalizations give \((t-1)/t\) and \(1/(1-t)\); multiplying these cyclic expressions yields the first identity, while eliminating the third gives the second. The half-turn symmetries about common perpendiculars of opposite edges exchange opposite edges and force the angle and parameter assertions.
\end{proof}

\begin{proposition}[Ideal-edge gluing equations]
\label{thurston-ch04-ideal-edge-gluing}
\group{hyperbolic geometry}
\source{thurston:page-0059,thurston:page-0060,thurston:page-0061}
\uses{thurston-ch04-shape-identities}
For tetrahedra glued around an edge with parameters \(z(e_1),\ldots,z(e_k)\), the hyperbolic structure extends across the edge precisely when
\[
 z(e_1)\cdots z(e_k)=1
\]
in \(\widetilde{\C^*}\), equivalently when the product is one and
\[
 \arg z(e_1)+\cdots+\arg z(e_k)=2\pi,\qquad 0<\arg z(e_i)\leq\pi.
\]
\notready
\end{proposition}

\begin{proof}
\sketch The face pairings already produce a hyperbolic structure off the one-skeleton. Around an edge, the angle equation removes rotational holonomy and the lifted product equation removes residual translation or dilation, exactly expressing extension over the edge.
\end{proof}

\begin{proposition}[Figure-eight shape family]
\label{thurston-ch04-figure-eight-shape-family}
\group{hyperbolic geometry}
\source{thurston:page-0061,thurston:page-0062,thurston:page-0063}
\uses{thurston-ch04-ideal-edge-gluing}
For the two ideal tetrahedra giving the figure-eight knot complement, with \(z=z_1\), \(w=w_1\), the independent gluing equation is
\[
 z(z-1)w(w-1)=1.
\]
Thus
\[
 z=\frac{1\pm\sqrt{1+4/(w(w-1))}}2.
\]
Set
\[
 \Delta(w)=1+\frac{4}{w(w-1)},\qquad
 R=\{w\in\C:\operatorname{Im}w>0\text{ and }\Delta(w)\notin[0,\infty)\}.
\]
For every \(w\in R\), exactly one choice of the square root gives \(\operatorname{Im}z>0\); these are precisely the solutions in which both tetrahedra are nondegenerate and positively oriented.  The slit in the boundary of \(R\) is the ray
\[
 r=\left\{\frac12+it:t\geq\frac{\sqrt{15}}2\right\}.
\]
At its endpoint the discriminant vanishes, and along the rest of \(r\) it is positive real, so no point of \(r\) gives two nondegenerate positively oriented tetrahedra.  The complete solution is \(z=w=\sqrt[3]{-1}=\frac12+\frac{\sqrt3}{2}i\).
\notready
\end{proposition}

\begin{proposition}[Completeness of the original solution]
\label{thurston-ch04-figure-eight-completeness}
\group{hyperbolic geometry}
\source{thurston:page-0064,thurston:page-0065}
\uses{thurston-ch04-figure-eight-shape-family}
For the peripheral generators \(x,y\), the holonomy derivatives are
\[
 H'(x)=\left(\frac zw\right)^2,\qquad H'(y)=w(1-z).
\]
Completeness forces both derivatives to be one, hence \(z=w\), and the gluing equation then forces the unique solution \(z=w=\sqrt[3]{-1}\).
\notready
\end{proposition}

\begin{proof}
\sketch Compute vector ratios in the cusp triangulation and multiply them around the two generators. Setting both similarities equal to the identity gives \(z=w\); the sign of \(z(z-1)\) selects \(z(z-1)=-1\), yielding the stated root.
\end{proof}

\section{Completion and generalized surgery}

\begin{proposition}[Completion at an ideal vertex]
\label{thurston-ch04-ideal-vertex-completion}
\group{hyperbolic geometry}
\source{thurston:page-0065,thurston:page-0066,thurston:page-0067}
Let \(\mathcal N(v)\) be a deleted neighborhood of an ideal vertex whose developing image is a solid cone with its axis removed. Its metric completion is obtained by adjoining the axis. If the modulus group \(\{|H'(\alpha)|\}\) is dense in \(\R_+\), the completion adds one point. If it is cyclic, the completion is \(\widetilde{\mathcal N}(v)/H\) and adds a circle of length \(\log|H'(\alpha_2)|\), where \(\alpha_2\) generates the modulus image. A transverse section is a cone of angle \(\arg\widetilde H'(\alpha_2)\); it is nonsingular exactly when this angle is \(2\pi\).
\notready
\end{proposition}

\begin{definition}[Generalized Dehn surgery invariant]
\label{thurston-ch04-generalized-surgery-invariant}
\group{hyperbolic geometry}
\source{thurston:page-0067,thurston:page-0068}
Let \(M\) have torus boundary components \(P_i\) with bases \((a_i,b_i)\). Write \(M_{(\alpha_i,\beta_i)}\) for filling in which the meridian maps to \(\alpha_i a_i+\beta_i b_i\), and write \(\infty\) for an unfilled component. For a hyperbolic completion, the generalized invariant \((\alpha_i,\beta_i)\) is the primitive-up-to-sign solution of
\[
 \alpha_i\widetilde H(a_i)+\beta_i\widetilde H(b_i)=\text{rotation by }\pm2\pi,
\]
with value \(\infty\) when the end is complete. Nonprimitive integral solutions give cone angles \(2\pi/n_i\), and the corresponding branched cover of index \(n_i\) is nonsingular.
\notready
\end{definition}

\section{Figure-eight surgery and degeneration}

\begin{proposition}[Figure-eight surgery parameter map]
\label{thurston-ch04-surgery-parameter-map}
\group{hyperbolic geometry}
\source{thurston:page-0069,thurston:page-0070,thurston:page-0071}
\uses{thurston-ch04-generalized-surgery-invariant}
Let \(\widehat{\R}^2=\R^2\cup\{\infty\}\) be the one-point compactification.  On the admissible shape region \(R\), the generalized surgery invariant first defines a continuous quotient-valued map
\[
 d:R\longrightarrow \widehat{\R}^2/\{v\sim-v\},
 \qquad d(w)=[(\mu,\lambda)]
\]
away from the complete shape \(w_0=\sqrt[3]{-1}\), while \(d(w_0)=\infty\); the antipodal relation is imposed on \(\R^2\) and fixes \(\infty\).  The orientation of the singular axis determined by the way the tetrahedral corners wrap around it chooses a lift \(\widetilde d:R\to\widehat{\R}^2\), with \(\widetilde d(w_0)=\infty\).  Away from \(w_0\), this lift is specified by \(\widetilde d(w)=(\mu,\lambda)\) and
\[
 \mu\widetilde H(m)+\lambda\widetilde H(l)=\text{rotation by }+2\pi,
\]
with \(m=y\), \(l=x+2y\), and
\[
 H(m)=w(1-z),\qquad H(l)=z^2(1-z)^2.
\]
The lift extends continuously to send the compactified boundary interval from \(1\) to \(+\infty\) onto the segment from \((-4,1)\) to \((4,1)\), and the interval from \(-\infty\) to \(0\) onto the segment from \((4,-1)\) to \((-4,-1)\).  If \(\tau\) is the involution induced by exchanging the shape parameters \(z\) and \(w\), and \(\sigma(w)=1/(1-w)\), then
\[
 \widetilde d(\tau w)=-\widetilde d(w),\qquad
 \widetilde d(\sigma w)=(\mu,-\lambda)
 \quad\text{when }\widetilde d(w)=(\mu,\lambda).
\]
These segment formulas do not describe the boundary near \(r\) or \(\tau(r)\).  Those boundary loci are not genuine boundaries of the family of hyperbolic structures: just beyond a crossing of \(r\), the \(z\)-tetrahedron is negatively oriented, but for a neighborhood past the crossing it can be cut into pieces and those pieces subtracted from the positively oriented \(w\)-tetrahedron.  The remaining polyhedron, with the induced face identifications, continues the hyperbolic structure.  Thus near \(r\) and \(\tau(r)\) this cut-and-subtracted-polyhedron continuation replaces the two-positively-oriented-tetrahedra description, and only qualitative boundary behavior is asserted.
\notready
\end{proposition}

\begin{theorem}[Hyperbolic Dehn surgery on the figure-eight knot]
\label{thurston-ch04-hyperbolic-dehn-surgery-theorem}
\dcref{Theorem 4.7}
\group{hyperbolic geometry}
\source{thurston:page-0071,thurston:page-0072}
\uses{thurston-ch04-surgery-parameter-map}
Every Dehn filling of the figure-eight knot complement has a hyperbolic structure except, up to simultaneous sign, the six slopes
\[
 (1,0),(0,1),(1,1),(2,1),(3,1),(4,1).
\]
\notready
\end{theorem}

\begin{definition}[Hyperbolic foliation]
\label{thurston-ch04-hyperbolic-foliation}
\group{foliations}
\source{thurston:page-0072,thurston:page-0073,thurston:page-0074}
A codimension-\(k\) foliation is a \(\G\)-structure locally modelled on \(\R^{n-k}\times\R^k\) with coordinate changes preserving horizontal \((n-k)\)-planes. Given a pseudogroup \(\H\) on a \(k\)-manifold \(N\), an \(\H\)-foliation has transverse changes in \(\H\). For hyperbolic isometries this is a codimension-\(k\) hyperbolic foliation; its developing map \(D:\widetilde M\to N\) is a submersion and its holonomy satisfies \(D\circ T_\alpha=H(\alpha)\circ D\).
\notready
\end{definition}

\begin{proposition}[Completeness criterion for foliations]
\label{thurston-ch04-foliation-developing-fibration}
\dcref{Proposition 4.8.1}
\group{foliations}
\source{thurston:page-0074,thurston:page-0075}
If the transverse group \(J\) is transitive and its isotropy groups are compact, the developing map of a differentiable \(\H\)-foliation on a closed manifold is a fibration.
\notready
\end{proposition}

\begin{proof}
\sketch Choose a transverse plane field and invariant transverse metric. Horizontal lifts have uniformly controlled length, hence extend globally by completeness of the universal cover. Lifting radial disks supplies local trivializations.
\end{proof}

\begin{theorem}[Classification of closed hyperbolic foliations]
\label{thurston-ch04-hyperbolic-foliation-classification}
\dcref{Theorem 4.9}
\group{foliations}
\source{thurston:page-0075,thurston:page-0081}
For a hyperbolic foliation of a closed three-manifold, either the holonomy group is discrete and the developing map descends to a Seifert fibration over \(\HH^2/H(\pi_1M)\), or the manifold fibers over \(S^1\) with torus fiber.
\notready
\end{theorem}

\begin{lemma}[Convex displacement function]
\label{thurston-ch04-convex-displacement}
\dcref{Lemma 4.9.1}
\group{foliations}
\source{thurston:page-0076,thurston:page-0077}
For orientation-preserving hyperbolic isometries \(g_1,g_2\), the function \(x\mapsto\dist(g_1x,g_2x)\) on \(\HH^2\) is convex.
\notready
\end{lemma}

\begin{proposition}[Classification of hyperbolic-plane isometries]
\label{thurston-ch04-hyperbolic-plane-isometries}
\dcref{Proposition 4.9.3}
\group{foliations}
\source{thurston:page-0078,thurston:page-0079}
Every nontrivial orientation-preserving isometry of \(\HH^2\) is exactly one of elliptic (one interior fixed point), hyperbolic (one invariant geodesic), or parabolic (one boundary fixed point). The centralizer of each such isometry is abelian.
\notready
\end{proposition}

\begin{proof}
\sketch Bisectors of successive segments \(x,gx,g^2x\) either meet, have a common perpendicular, or meet at infinity; these alternatives give the three types and their uniqueness. The invariant-point or invariant-axis description makes the centralizer abelian.
\end{proof}

\begin{proposition}[Torus bundle in the nondiscrete case]
\label{thurston-ch04-nondiscrete-foliation-torus-bundle}
\group{foliations}
\source{thurston:page-0076,thurston:page-0077,thurston:page-0080,thurston:page-0081}
\uses{thurston-ch04-convex-displacement,thurston-ch04-hyperbolic-plane-isometries}
If the holonomy of a closed hyperbolic foliation is nondiscrete, its small elements contain a least positive element in the induced ordering; this element is parabolic. The invariant point at infinity yields a nonsingular closed one-form whose periods are discrete, so integration fibers the manifold over \(S^1\) with torus fibers. The same conclusion holds after passing to, and descending from, the double cover which orients the foliation.
\notready
\end{proposition}

\begin{example}[Hyperbolic torus bundles]
\label{thurston-ch04-hyperbolic-torus-bundles}
\group{foliations}
\source{thurston:page-0081}
If \(A\in SL(2,\Z)\) is hyperbolic, its mapping torus has two codimension-two hyperbolic foliations from the eigendirections. For \(A=\begin{bmatrix}2&1\\1&1\end{bmatrix}\), the mapping torus is a figure-eight Dehn filling with slope \((D,\pm1)\). All codimension-two hyperbolic foliations with nonclosed leaves arise this way.
\notready
\end{example}

\section{Incompressible surfaces}

\begin{definition}[Incompressibility and Haken manifold]
\label{thurston-ch04-incompressibility}
\group{3-manifold topology}
\source{thurston:page-0082,thurston:page-0083}
A surface \(S\subset M\), excluding the elementary cases stated in the source, is incompressible if every loop on \(S\) bounding a disk in \(M-S\) bounds a disk in \(S\), and is \(\partial\)-incompressible if every boundary arc homotopic into \(\partial M\) is homotopic in \(S\) into \(\partial S\). An oriented irreducible compact three-manifold containing such a surface is sufficiently large, or Haken.
\notready
\end{definition}

\begin{proposition}[The six essential surfaces]
\label{thurston-ch04-figure-eight-essential-surfaces}
\group{3-manifold topology}
\source{thurston:page-0083,thurston:page-0084}
For the figure-eight knot exterior \(M\), the six distinguished connected incompressible and \(\partial\)-incompressible surfaces are: the boundary-parallel torus \(S_1\); a punctured torus Seifert surface \(S_2\) with boundary slope \((0,\pm1)\); two punctured Klein bottles \(S_3,S_5\); and the boundaries \(S_4,S_6\) of their tubular neighborhoods, each a twice-punctured torus with slope \((\pm4,1)\).
\notready
\end{proposition}

\begin{theorem}[Classification of essential surfaces]
\label{thurston-ch04-essential-surface-classification}
\dcref{Theorem 4.11}
\group{3-manifold topology}
\source{thurston:page-0083,thurston:page-0084,thurston:page-0085,thurston:page-0086,thurston:page-0087,thurston:page-0088,thurston:page-0089,thurston:page-0090,thurston:page-0091,thurston:page-0092,thurston:page-0093,thurston:page-0094}
\uses{thurston-ch04-incompressibility,thurston-ch04-figure-eight-essential-surfaces}
Every connected incompressible and \(\partial\)-incompressible surface in the figure-eight knot exterior is isotopic to one of \(S_1,\ldots,S_6\).
\notready
\end{theorem}

\begin{proof}
\sketch Use the two-hexagon spine of the exterior and isotope the surface to remove circular intersections, return arcs, and unnecessary corner connections while minimizing intersections with the one-skeleton. If all four edge counts are positive, the corner equations force the unique pattern yielding parallel punctured-torus fibers. If an edge count vanishes, minimality leaves only the punctured-torus, twice-punctured-torus, or punctured-Klein-bottle patterns. In the complementary \(T^2\times I\) region, incompressibility reduces the surface to boundary-parallel pieces or vertical annuli; connectedness gives exactly the five non-boundary cases \(S_2,\ldots,S_6\).
\end{proof}

\begin{corollary}[Irreducibility and Haken fillings]
\label{thurston-ch04-haken-surgery-corollary}
\group{3-manifold topology}
\source{thurston:page-0084,thurston:page-0085}
Every Dehn surgery filling \(M_{(m,l)}\) of the figure-eight knot exterior is irreducible, and it is Haken exactly for slopes \((0,\pm1)\) and \((\pm4,\pm1)\). The other hyperbolic fillings are irreducible, non-Haken, and not Seifert fiber spaces.
\notready
\end{corollary}

\begin{proof}
\sketch Put a sphere or incompressible surface in minimal position with the filling solid torus. Compressing or boundary-compressing the exterior intersection decreases core intersections until it is one of the classified six surfaces. None can occur as a sphere, and only the listed boundary slopes cap the essential surfaces, proving irreducibility and the Haken criterion.
\end{proof}
